% GENERATED FILE. Edit canonical lessons, then run npm run generate:books.
% canonical-source-hash: fnv1a64:9587be207122f3c9
% canonical-lessons: BN-C488-shoishob, BN-C488-koishor, BN-C488-joubon, BN-C488-bardhokyo, BN-C488-bere-otha

\chapter{Growing Up and Growing Old}
\label{ch:bn-growing-up-and-growing-old}
\hlchaptermodality{488}

\begin{chapteropening}
\textbf{By the end of this chapter:} \emph{I can talk about childhood, youth and old age.}

Complete the last lesson of chapter 488: I can talk about childhood, youth and old age.
\end{chapteropening}

\section[\texorpdfstring{\bn{শৈশব} (shoishôb)}{শৈশব (shoishôb)}]{\bn{শৈশব} (shoishôb) — childhood}

\label{lesson:BN-C488-shoishob}



\begin{quote}
Before the new one: say the Bengali for to guard protectively, then the Bengali for to stroke.
\end{quote}

\subsection*{You'll want to know: \bn{শৈশব}}
\textbf{\bn{শৈশব}} — \emph{shoishôb} — \textquotedblleft{}childhood\textquotedblright{}.

\begin{quote}
\bn{দাদুর} \bn{শৈশব} \bn{গ্রামে} \bn{কাটল।}
\end{quote}

\begin{quote}
\emph{dādur shoishôb grāme kāṭlô.}
\end{quote}

\begin{quote}
Grandfather spent his childhood in the village.
\end{quote}

\begin{cousinweb}
\textbf{\bn{শৈশব}} (\emph{shoishôb}) is a tatsama, Sanskrit \emph{śaiśava}, \textquotedblleft{}childhood\textquotedblright{} (Monier-Williams), from \emph{śiśu}, the \bn{শিশু} (taught) of Bengali.
\end{cousinweb}

\begin{grammarlens}[title={What you've built}]
One more word for describing what you lived and felt.
\end{grammarlens}

\subsection*{Your turn}
\begin{itemize}
\raggedright
  \item \emph{Say it:} \emph{shoishôb}
  \item \emph{Say it:} \emph{shoishôb}, once more
  \item \emph{Say it:} \emph{dādur shoishôb grāme kāṭlô.} — Grandfather spent his childhood in the village.
  \item \emph{Recall:} say the Bengali for to guard protectively, then the Bengali for to stroke, then say \emph{shoishôb} again
\end{itemize}

\begin{tcolorbox}[breakable,colback=teal!4,colframe=teal!35!black,title={Before you move on}]
What does \bn{শৈশব} mean? (\textquotedblleft{}Childhood\textquotedblright{}.) Say it once more.
\end{tcolorbox}
\section[\texorpdfstring{\bn{কৈশোর} (koishor)}{কৈশোর (koishor)}]{\bn{কৈশোর} (koishor) — adolescence; the early teens}

\label{lesson:BN-C488-koishor}



\begin{quote}
Before the new one: say the Bengali for to stroke, then the Bengali for childhood.
\end{quote}

\subsection*{You'll want to know: \bn{কৈশোর}}
\textbf{\bn{কৈশোর}} — \emph{koishor} — \textquotedblleft{}adolescence; the early teens\textquotedblright{}.

\begin{quote}
\bn{কৈশোরে} \bn{অরুণ} \bn{ফুটবল} \bn{খেলত।}
\end{quote}

\begin{quote}
\emph{koishore ôrun phuṭbôl kheltô.}
\end{quote}

\begin{quote}
As a teenager Arun used to play football.
\end{quote}

\begin{cousinweb}
\textbf{\bn{কৈশোর}} (\emph{koishor}) is a tatsama, Sanskrit \emph{kaiśora}, \textquotedblleft{}youth, boyhood from ten to fifteen\textquotedblright{} (Monier-Williams).
\end{cousinweb}

\begin{grammarlens}[title={What you've built}]
One more word for describing what you lived and felt.
\end{grammarlens}

\subsection*{Your turn}
\begin{itemize}
\raggedright
  \item \emph{Say it:} \emph{koishor}
  \item \emph{Say it:} \emph{koishor}, once more
  \item \emph{Say it:} \emph{koishore ôrun phuṭbôl kheltô.} — As a teenager Arun used to play football.
  \item \emph{Recall:} say the Bengali for to stroke, then the Bengali for childhood, then say \emph{koishor} again
\end{itemize}

\begin{tcolorbox}[breakable,colback=teal!4,colframe=teal!35!black,title={Before you move on}]
What does \bn{কৈশোর} mean? (\textquotedblleft{}Adolescence; the early teens\textquotedblright{}.) Say it once more.
\end{tcolorbox}
\section[\texorpdfstring{\bn{যৌবন} (joubôn)}{যৌবন (joubôn)}]{\bn{যৌবন} (joubôn) — youth; young adulthood}

\label{lesson:BN-C488-joubon}



\begin{quote}
Before the new one: say the Bengali for childhood, then the Bengali for adolescence.
\end{quote}

\subsection*{You'll want to know: \bn{যৌবন}}
\textbf{\bn{যৌবন}} — \emph{joubôn} — \textquotedblleft{}youth; young adulthood\textquotedblright{}.

\begin{quote}
\bn{যৌবনে} \bn{দাদু} \bn{সেনাবাহিনীতে} \bn{ছিলেন।}
\end{quote}

\begin{quote}
\emph{joubône dādu senābāhinīte chhilen.}
\end{quote}

\begin{quote}
In his youth my grandfather was in the army.
\end{quote}

\begin{cousinweb}
\textbf{\bn{যৌবন}} (\emph{joubôn}) is a tatsama, Sanskrit \emph{yauvana}, \textquotedblleft{}youth, youthfulness\textquotedblright{} (Monier-Williams).
\end{cousinweb}

\begin{grammarlens}[title={What you've built}]
One more word for describing what you lived and felt.
\end{grammarlens}

\subsection*{Your turn}
\begin{itemize}
\raggedright
  \item \emph{Say it:} \emph{joubôn}
  \item \emph{Say it:} \emph{joubôn}, once more
  \item \emph{Say it:} \emph{joubône dādu senābāhinīte chhilen.} — In his youth my grandfather was in the army.
  \item \emph{Recall:} say the Bengali for childhood, then the Bengali for adolescence, then say \emph{joubôn} again
\end{itemize}

\begin{tcolorbox}[breakable,colback=teal!4,colframe=teal!35!black,title={Before you move on}]
What does \bn{যৌবন} mean? (\textquotedblleft{}Youth; young adulthood\textquotedblright{}.) Say it once more.
\end{tcolorbox}
\section[\texorpdfstring{\bn{বার্ধক্য} (bārdhôkyô)}{বার্ধক্য (bārdhôkyô)}]{\bn{বার্ধক্য} (bārdhôkyô) — old age}

\label{lesson:BN-C488-bardhokyo}



\begin{quote}
Before the new one: say the Bengali for adolescence, then the Bengali for youth.
\end{quote}

\subsection*{You'll want to know: \bn{বার্ধক্য}}
\textbf{\bn{বার্ধক্য}} — \emph{bārdhôkyô} — \textquotedblleft{}old age\textquotedblright{}.

\begin{quote}
\bn{বার্ধক্যে} \bn{দাদু} \bn{একা} \bn{থাকেন।}
\end{quote}

\begin{quote}
\emph{bārdhôkye dādu ekā thāken.}
\end{quote}

\begin{quote}
In his old age my grandfather lives alone.
\end{quote}

\begin{cousinweb}
\textbf{\bn{বার্ধক্য}} (\emph{bārdhôkyô}) is a tatsama, Sanskrit \emph{vārddhakya}, \textquotedblleft{}old age\textquotedblright{} (Monier-Williams).
\end{cousinweb}

\begin{grammarlens}[title={What you've built}]
One more word for describing what you lived and felt.
\end{grammarlens}

\subsection*{Your turn}
\begin{itemize}
\raggedright
  \item \emph{Say it:} \emph{bārdhôkyô}
  \item \emph{Say it:} \emph{bārdhôkyô}, once more
  \item \emph{Say it:} \emph{bārdhôkye dādu ekā thāken.} — In his old age my grandfather lives alone.
  \item \emph{Recall:} say the Bengali for adolescence, then the Bengali for youth, then say \emph{bārdhôkyô} again
\end{itemize}

\begin{tcolorbox}[breakable,colback=teal!4,colframe=teal!35!black,title={Before you move on}]
What does \bn{বার্ধক্য} mean? (\textquotedblleft{}Old age\textquotedblright{}.) Say it once more.
\end{tcolorbox}
\section[\texorpdfstring{\bn{বেড়ে} \bn{ওঠা} (beṛe oṭhā)}{বেড়ে ওঠা (beṛe oṭhā)}]{\bn{বেড়ে} \bn{ওঠা} (beṛe oṭhā) — to grow up}

\label{lesson:BN-C488-bere-otha}



\begin{quote}
Before the new one: say the Bengali for youth, then the Bengali for old age.
\end{quote}

\subsection*{You'll want to know: \bn{বেড়ে} \bn{ওঠা}}
\textbf{\bn{বেড়ে} \bn{ওঠা}} — \emph{beṛe oṭhā} — \textquotedblleft{}to grow up\textquotedblright{}.

\begin{quote}
\bn{অরুণ} \bn{গ্রামে} \bn{দাদুর} \bn{কাছে} \bn{বেড়ে} \bn{উঠল।}
\end{quote}

\begin{quote}
\emph{ôrun grāme dādur kāchhe beṛe uṭhlô.}
\end{quote}

\begin{quote}
Arun grew up with his grandfather in the village.
\end{quote}

Past (he/she …): \bn{তিনি} \bn{শহরে} \bn{বেড়ে} \bn{উঠলেন}, \emph{tini shôhôre beṛe uṭhlen}.

\begin{cousinweb}
\textbf{\bn{বেড়ে} \bn{ওঠা}} (\emph{beṛe oṭhā}) is \emph{beṛe}, the -\bn{এ} form of \emph{bāṛā}, \textquotedblleft{}to grow\textquotedblright{}, which CDIAL lists for Bengali under Sanskrit \emph{vardhate}, \textquotedblleft{}grows\textquotedblright{} (CDIAL 11376), with \emph{oṭhā}, \textquotedblleft{}to rise\textquotedblright{} (CDIAL 1793).
\end{cousinweb}

\begin{grammarlens}[title={What you've built}]
One more word for describing what you lived and felt.
\end{grammarlens}

\subsection*{Your turn}
\begin{itemize}
\raggedright
  \item \emph{Say it:} \emph{beṛe oṭhā}
  \item \emph{Say it:} \emph{beṛe oṭhā}, once more
  \item \emph{Say it:} \emph{ôrun grāme dādur kāchhe beṛe uṭhlô.} — Arun grew up with his grandfather in the village.
  \item \emph{Recall:} say the Bengali for youth, then the Bengali for old age, then say \emph{beṛe oṭhā} again
\end{itemize}

\begin{tcolorbox}[breakable,colback=teal!4,colframe=teal!35!black,title={Before you move on}]
What does \bn{বেড়ে} \bn{ওঠা} mean? (\textquotedblleft{}To grow up\textquotedblright{}.) Say it once more.
\end{tcolorbox}
